\documentclass[12pt]{article}

\usepackage{amsmath}
\usepackage{parskip}
\usepackage{amssymb}
\usepackage{geometry}
\usepackage{hyperref}
\usepackage{booktabs}
\usepackage{float}
\usepackage{graphicx}
\usepackage{longtable}
\usepackage{url}
\usepackage{natbib}
\usepackage{pdfpages}
\hypersetup{colorlinks=true, linkcolor=blue}

\geometry{a4paper, margin=1in}

\date{May 1 2026}

\title{Nichols Radial Injection Model XI: The Sequestration of r-Process Elements and Widowed Baryonic Matter}
\author{Lawrence William Nichols}

\begin{document}

\maketitle

\begin{abstract}
Although this remains a hypothesis, it is quite possible that r-process elements (gold, platinum, uranium, thorium, rare earths, and other nuclei with $Z > 28$) are rare in the observable universe today. The majority appear to have been produced in the very early universe and subsequently sequestered into ``widowed'' baryonic matter --- out-of-time, Lorentz-synchronised material ejected via extreme Hawking decoupling at black-hole accelerators. This material remains electromagnetically decoupled from our local cosmic time frame while still contributing gravitationally, providing a natural explanation for the extreme deuterium enrichment observed in interstellar object 3I/ATLAS as the coupled light-fraction fossil of an ancient widowed r-process core.
\end{abstract}

\section{Observational Evidence}

The systematic absence of r-process elements is consistently observed across multiple independent datasets:

\begin{itemize}
    \item Lunar samples (Apollo 11--17, Luna, Chang'e-5 missions) and in-situ Mars rover analyses (Curiosity APXS, Perseverance PIXL; Mars Odyssey GRS orbital data) exhibit the identical cutoff: major-element inventories are strictly light-to-iron-peak ($Z \leq 28$), with r-process nuclei (REE, Th, U) present only at trace/diluted levels or confined to minor components such as lunar KREEP (Th $\sim$1--15\,ppm, U $\sim$0.2--4\,ppm; REE enriched locally but never major). These pristine, human-uncontaminated materials reinforce that the sequestration pattern is primordial and solar-system-wide rather than an artifact of terrestrial processing or late veneer delivery.
    
    \item Black-hole accretion environments and relativistic jets (across stellar-mass to supermassive regimes; 25+ objects) show only light-to-iron-peak elements ($Z \leq 28$: H, He, C, N, O, Ne, Na, Mg, Si, S, Ca, Cr, Fe, Ni) plus common molecules (H$_2$, H$_2$O, CO, HCN, HCO$^+$, CS, SiO, etc.).
    
    \item Primitive meteorites (iron, pallasites, chondrites, mesosiderites; 50 representative specimens) exhibit the same chemical cutoff at $Z = 28$, with dominant Fe--Ni metal, olivine, pyroxene, and CAIs (Ca--Al); heavies are diluted or absent from major-element inventories.
    
    \item Most GRB afterglows and associated supernovae lack clear r-process signatures, despite being among the most violent and neutron-rich events known.
\end{itemize}

\begin{table}[htbp]
\centering
\caption{Lunar Samples \& Mars Rover/Orbital Data}
\begin{tabular}{p{8.0cm}c}
\toprule
\textbf{Light-to-Iron-Peak ($Z \leq 28$) + Minerals} & \textbf{r-Process ($Z > 28$)} \\
\midrule
O, Si, Mg, Fe, Ca, Al, Ti (dominant in basalts, anorthosites, regolith) \\
KREEP component (minor); olivine, pyroxene, plagioclase 
& NULL in major inventories \\
(REE, Th, U at ppm traces only; Au/Pt highly depleted) \\
\bottomrule
\end{tabular}
\end{table}

\begin{table}[htbp]
\centering
\caption{Black-Hole Accretion \& Relativistic Jets (25+ objects)}
\begin{tabular}{p{8.0cm}c}
\toprule
\textbf{Light-to-Iron-Peak ($Z \leq 28$) + Molecules} & \textbf{r-Process ($Z > 28$)} \\
\midrule
H, He, C, N, O, Ne, Na, Mg, Si, S, Ca, Cr, Fe, Ni \\
H$_2$, H$_2$O, CO, HCN, HCO$^+$, CS, SiO 
& NULL \\
\bottomrule
\end{tabular}
\end{table}

\begin{table}[htbp]
\centering
\caption{Primitive Meteorites (iron, pallasites, chondrites, mesosiderites; 50 representative specimens)}
\begin{tabular}{p{8.0cm}c}
\toprule
\textbf{Light-to-Iron-Peak ($Z \leq 28$) + Molecules} & \textbf{r-Process ($Z > 28$)} \\
\midrule
Fe--Ni dominant + Mg, Si, O, Ca, Al, C, H, N \\
(olivine, pyroxene, CAIs; heavies diluted)
& NULL \\
\bottomrule
\end{tabular}
\end{table}

This pattern suggests that the bulk of r-process production occurred early and that most of the resulting material was locked into widowed baryonic cores rather than efficiently mixed into the later interstellar medium.

\section{Implications for Cosmic Chronology}

Only a small fraction of r-process material escaped sequestration and was delivered to the solar system. Gold-197, being stable and monoisotopic, cannot be directly radiometrically dated. However, the incorporation of gold into the solar system is tightly constrained by Calcium-Aluminium-rich Inclusions (CAIs) in primitive meteorites at approximately \textbf{4.567 billion years ago}, with much of Earth's accessible gold arriving via the late veneer around \textbf{3.9 billion years ago}. Even where (U--Th)/He dating of native gold and associated inclusions records mineralization ages of order 100\,Myr, these ages represent only the final re-synchronization and hydrothermal crystallization phase of the escaped light fraction. Within the Widowed Matter Hypothesis the material must therefore have remained Lorentz-desynchronized for $\sim$13.7\,Gyr of the 13.8\,Gyr cosmic timeline. Successful measurement of the true nucleosynthesis age locked inside desynchronized heavy r-process cores (via a future method sensitive to desynchronized proper time) would thus directly constrain the actual age of the universe.

Some short-lived radionuclides exhibit half-lives measured in years rather than millions of years, yet all are observed to have decayed on comparable short timescales in early Solar System materials. Within the Widowed Matter Hypothesis, these half-lives do not reflect the full cosmic age since nucleosynthesis. Instead, they record the catch-up rate at which partially desynchronized material re-synchronizes with our local cosmic time frame after ejection at the kinematic wall (at the kinematic wall, \(v \approx 0.9c\), \(\gamma \approx 2.29\)).

If a method to probe the ``age'' of the widowed (out-of-time) fraction ever becomes possible, it could provide a direct test of this hypothesis and help determine the true effective age of the universe.

The light fraction (cyan spheres) gradually returns to our clock while the heavy r-process core remains Lorentz-desynchronized, naturally explaining the remarkably brief apparent intervals recorded in CAIs and primitive meteorites.

\section{Testable Predictions}

\begin{itemize}
    \item Interstellar objects such as 3I/ATLAS may contain dense, partially re-synced widowed r-process cores, producing anomalous outgassing, non-gravitational acceleration, and unusual volatile chemistry without prominent heavy-element spectral lines.
    \item Future high-precision microlensing or dynamical studies of dark compact objects may reveal indirect signatures of r-process enrichment.
\end{itemize}

If these predictions hold, the sequestration of r-process elements into widowed baryonic matter offers a unified mechanical explanation for both the ``missing heavies'' in observable environments and the composition of dark matter.

\section{Video Catalog}
A video overview of the full RIM series is available on the author's \href{https://www.youtube.com/@Nichols_Thought_Experiments}{Nichols Radial Injection Model (RIM)} YouTube channel.

\begin{thebibliography}{99}

\bibitem{rimI-X}
L.\,W.\,Nichols.
Nichols Radial Injection Model (RIM) Papers I--X [Self-published], January to May 2026.
\newline
\url{https://rxiverse.org/author/lawrence_william_nichols}

\bibitem{taylor2014}
Taylor, G.\,J. (2014).
Lunar bulk chemical composition: a post-Gravity Recovery and Interior Laboratory re-assessment.
\textit{Philosophical Transactions of the Royal Society A}, 372(2024), 20130242.
https://doi.org/10.1098/rsta.2013.0242

\bibitem{warren1979}
Warren, P.\,H. (1979).
The origin of KREEP.
\textit{Reviews of Geophysics}, 17(1), 73--88.

\bibitem{laul1973}
Laul, J.\,C., et al. (1973).
Chemical composition of Apollo 15, 16, and 17 samples.
\textit{Proceedings of the Lunar Science Conference}, 4, 1349--1367.

\bibitem{boynton2007}
Boynton, W.\,V., et al. (2007).
Elemental abundances determined via the Mars Odyssey GRS.
In \textit{The Martian Surface: Composition, Mineralogy, and Physical Properties} (pp. 105--124). Cambridge University Press.

\bibitem{gellert2014}
Gellert, R., et al. (2006--present series).
Mars Exploration Rovers and MSL/Curiosity APXS results (multiple papers).
Key reference: Bruckner, J., et al. (2014).
Mars surface chemistry: APXS results.
In \textit{The Martian Surface}.

\bibitem{allwood2021}
Allwood, A.\,C., et al. (2021).
The PIXL Instrument on the Mars 2020 Perseverance Rover.
arXiv:2103.07001 (and subsequent Science papers 2022--2026).

\end{thebibliography}

\newpage

\section{Engineering Studies and Visual Data}
This section includes relevant visual diagrams.
\includepdf[pages=-]{r-process.pdf}

\end{document}